\section{Marking and Refinement}
\label{sec:marking-refinement}

The localised DWR contributions defined in \cref{sec:space-time-local-indicator-construction} are associated with space--time cells \(Q_{K,n}=K\times I_n\). Applying a standard marking strategy like D\"orfler marking \cite{Dorfler1996} would lead to an unstructured mesh in space-time. However, we aim to retain the tensor product structure of our space–time mesh, while allowing the spatial mesh to vary from slab to slab, as this allows for standard timestepping. We therefore need to devise a specialised marking strategy. 

We separate the our refinement decisions into spatial and temporal components. With \(\eta_{K,n}\) denoting the signed local DWR contribution associated with \(Q_{K,n}\), we denote their absolute values
\begin{equation}
  s_{K,n}\coloneqq |\eta_{K,n}|,
  \qquad
  S\coloneqq
  \sum_{n=1}^{N}\sum_{K\in\mathcal T_n}s_{K,n}.
  \label{eq:marking-activity}
\end{equation}

We prescribe the D\"orfler marking parameter \(\theta\in(0,1]\), denoting the fraction of the total sum that should be reached by the refinement. The global space--time marking set \(\mathcal M_{\mathrm{st}}\) is then chosen as the smallest set of pairs \((K,n)\) satisfying
\begin{equation}
  \sum_{(K,n)\in\mathcal M_{\mathrm{st}}}s_{K,n}
  \ge \theta S.
  \label{eq:space-time-dorfler}
\end{equation}

For each time slab \(I_n\), we collect the marked space--time cells whose second index is \(n\). This gives the set of spatial cells \(\mathcal M_n\) to be refined on that slab:
\begin{equation}
  \mathcal M_n
  \coloneqq
  \left\{
    K\in\mathcal T_n:
    (K,n)\in\mathcal M_{\mathrm{st}}
  \right\}.
  \label{eq:slabwise-marked-cells}
\end{equation}
The spatial mesh associated with \(I_n\) is then refined only in these marked cells:
\begin{equation}
  \mathcal T_n^{\mathrm{new}}
  \coloneqq
  \operatorname{refine}(\mathcal T_n,\mathcal M_n).
  \label{eq:slabwise-spatial-refinement}
\end{equation}

The temporal refinement decision is made independently at slab level. A slab is bisected when the marked cells occupy at least a prescribed fraction \(\varphi_t\) of its spatial mesh:
\begin{equation}
  \frac{|\mathcal M_n|}{|\mathcal T_n|}
  \ge \varphi_t.
  \label{eq:within-slab-time-trigger}
\end{equation}
If \(I_n\) is bisected, both temporal children inherit the refined spatial mesh \(\mathcal T_n^{\mathrm{new}}\). 

The complete procedure therefore consists of one global space--time bulk selection, followed by slab-local spatial refinement and a separate slabwise decision on temporal bisection.



\section{Automated Adaptive Algorithm}

The complete solve--estimate--mark--refine cycle is summarised in \cref{alg:complete-adaptive-cycle}.

\begin{algorithm}[htbp]
\caption{Goal-oriented slabwise space--time adaptive cycle}
\label{alg:complete-adaptive-cycle}
\begin{algorithmic}[1]
\Require PDE, goal functional \(J\), initial slab meshes, estimator type, bulk parameter \(\theta\), temporal threshold \(\varphi_t\)

\State Initialise the current space--time discretisation

\While{the stopping criterion is not satisfied}
  \State Construct the current and enriched spatial spaces, together with the transfer operators \(P_{n-1}\) and \(P_{n-1}^{\ast}\) defined in \eqref{eq:forward-slab-transfer} and \eqref{eq:mass-adjoint-transfer-definition}.
  
  \State Solve the primal problem using the dG-in-time formulation \eqref{eq:global-dg-semilinear-form} to algebraic
  tolerance \(\mathrm{TOL}_{\mathrm{alg}}\).

  \If{the nonlinear estimator is selected}
    \State Solve the enriched primal problem according to \eqref{eq:enriched-primal}.
    \State Solve the current and enriched adjoint problems according to \eqref{eq:discrete-nonlinear-adjoint} and \eqref{eq:enriched-adjoint}.
    \State Form \(e_u\) and \(e_z\) as defined in \eqref{eq:nonlinear-dwr-discrete-errors}.
    \State Assemble \(\eta_{\mathrm{global}}\) from \eqref{eq:practical-nonlinear-estimator}.
  \Else
  \State Solve the enriched adjoint problem defined by \eqref{eq:discrete-nonlinear-adjoint}. 
  \State Form \(z^\star=z_h^+-I_hz_h^+\) according to \eqref{eq:computable-dwr-interpolated-weight}.
    \State Assemble \(\eta_{\mathrm{global}}\) from \eqref{eq:computable-dwr-two-solves}.
  \EndIf
  
  \State Recover the local residual densities using \eqref{eq:space-time-volume-projection}--\eqref{eq:mixed-ridge-recovery}.
  
  \State Assemble the local entity contributions from \eqref{eq:complete-nonlinear-entity-indicator} and the space--time cell indicators \(\eta_{K,n}\) from \eqref{eq:complete-nonlinear-cell-indicator}.
  
  \If{\(|\eta_{\mathrm{global}}|\le\mathrm{TOL}\)}
    \State \Return the final discretisation and estimators.
  \EndIf
  
  \State Apply the space--time marking criterion \eqref{eq:space-time-dorfler} and the temporal criterion \eqref{eq:within-slab-time-trigger}.
  \State Refine the selected spatial meshes and bisect the selected temporal slabs.
\EndWhile
\end{algorithmic}
\end{algorithm}



\section{Code Implementation}
\label{sec:single-indicator-loop}

\subsection{Primal, Goal, and Adjoint Solves}
\label{sec:primal-goal-adjoint-solves}

The time-dependent problems are supplied as semidiscrete UFL residuals using Firedrake and Irksome. Firedrake provides the spatial finite-element spaces and variational solvers, while Irksome's \texttt{TimeStepper} makes the fully-discrete variational problems by applying the chosen discontinuous Galerkin time discretisation \cite{andrews2026automatedgalerkintimestepping}.

% \textcolor{blue}{Using pyadjoint and discard the current ufl-generated adjoint solver would be a better idea for automation. Should I rerun the experiments and so only pyadjoint would be included here? or maybe I'll list this as a future work that will be carried out soon}
The present implementation constructs the discrete adjoint explicitly from the UFL representation. In compact form, the discrete adjoint satisfies
\[
  \mathcal F_h'(U_h)^*Z_h=J_h'(U_h).
\]
UFL \texttt{derivative} constructs the linearisation of the primal residual, \texttt{adjoint} transposes the resulting bilinear form, and \texttt{action} inserts the current adjoint coefficient. For nonlinear problems, the resulting form is evaluated at the saved primal polynomial at the corresponding physical time. The common solver supplies the reverse-time mass term and solves the resulting linear problem slab by slab using the same Irksome dG machinery.

Terminal and running contributions to the goal functional are handled separately because they enter the reverse-time problem differently. The derivative of a terminal goal determines the final adjoint state, whose sensitivity is subsequently propagated through all preceding slabs. The derivative of a running goal instead acts as a source throughout the backward solve. Both derivatives are generated symbolically from the user-supplied UFL goal forms.

%A more fully automated alternative is provided by \texttt{firedrake-adjoint}, following the high-level discrete-adjoint methodology of \cite{Farrell_Ham_Funke_Rognes_2013} and implemented using Pyadjoint \cite{Mitusch2019}. It records the forward operations on a tape and traverses them in reverse to generate the corresponding discrete adjoint. This removes the need to write an explicit reverse-time solver.

The current implementation retains the complete forward trajectory in memory. For longer simulations, Firedrake's automated adjoint derivation supports checkpointing, which stores only selected forward states and recomputes intermediate states during the reverse traversal. This reduces memory consumption at the cost of additional forward computation. \cite{checkpointing}

% https://www.firedrakeproject.org/checkpointing.html#using-checkpointing-in-adjoint-simulations


\subsection{Residual Localisation and Refinement}

The primal residual is localised using a hierarchical bubble--cone recovery. Firedrake's scalar Bubble (\texttt{B}) and FacetBubble (\texttt{FB}) element families provide the auxiliary cell-interior and facet-supported test functions, respectively \cite{Firedrake}.
Temporal integration is performed using Gauss--Legendre quadrature, while the spatial forms are assembled by Firedrake using an explicitly selected quadrature degree.

In the two-dimensional examples, the initial geometry-conforming meshes and the marked-cell spatial refinements are generated using \texttt{Netgen} (cite).
\yue{Add Netgen}

% https://www.firedrakeproject.org/variational-problems.html


\subsection{Transfer on Nonmatching Slab Meshes}
\label{sec:slab-transfer-adjoint}

Following the refinement strategy in \cref{sec:marking-refinement}, consecutive time slabs may use different spatial meshes and finite-element spaces. The outgoing primal trace from \(I_{n-1}\) must therefore be transferred into \(V_n\). We denote the resulting linear interpolation operator by
\begin{equation}
  P_{n-1}:V_{n-1}\longrightarrow V_n
  \label{eq:forward-slab-transfer}
\end{equation}
and define the transferred dG interface defect by
\begin{equation}
  \delta_{n-1}(U_h)
  =
  U_h(t_{n-1}^{+})
  -
  P_{n-1}U_h(t_{n-1}^{-}).
  \label{eq:transferred-interface-defect-full}
\end{equation}

The reverse propagation is determined by the same discrete forward operator. Let \(m_n(\cdot,\cdot)\) denote the discrete mass pairing on \(V_n\), and let \( \mathbf{M_n}\) be its matrix representation. The adjoint transfer \(P_{n-1}^{\ast}:V_n\to V_{n-1}\) is defined by
\begin{equation}
  m_n(P_{n-1}w,z_n)
  =
  m_{n-1}(w,P_{n-1}^{\ast}z_n)
  \qquad
  \forall w\in V_{n-1},\quad z_n\in V_n.
  \label{eq:mass-adjoint-transfer-definition}
\end{equation}
\pef{I think the notation $M_n$ is new? Before it was $m$}
\yue{change to $m_n$}
In coefficient form,
\begin{equation}
  \mathbf P_{n-1}^{\ast} = \mathbf M_{n-1}^{-1} \mathbf P_{n-1}^{T} \mathbf M_n.
  \label{eq:mass-adjoint-transfer-matrix}
\end{equation}

In the implementation, Firedrake assembles \(P_{n-1}\) as a cross-mesh finite-element interpolation operator. The forward trace is obtained by applying this matrix, while the reverse transfer is evaluated by a target mass action, a transpose matrix action, and a source mass solve as stated in ~\eqref{eq:mass-adjoint-transfer-matrix}.

The construction above is formulated at the nodal interpolation and does not impose physical constraints. For applications involving conserved, or divergence-free quantities, however, the transfer \(P_{n-1}\) must be carefully chosen to preserve the relevant physical properties, see, e.g., \cite{non-flow}.

